\documentclass{article}
\usepackage[T1]{fontenc}
\usepackage{lmodern}
\usepackage{graphicx}
\usepackage{amsmath,amssymb}
\usepackage[a4paper,margin=2cm]{geometry}
\usepackage[absolute,overlay]{textpos}
\setlength{\TPHorizModule}{1mm}
\setlength{\TPVertModule}{1mm}
\usepackage[indonesian]{babel}
\usepackage{hyperref}
\setlength{\emergencystretch}{2em}
% unit: Halaman 12

\begin{document}

% === Halaman 12 (python/native) ===

Rinaldi Munir/II4021 Kriptografi

12

Medan Berhingga Fp

• Kelas medan berhingga yang penting adalah Fp ,  dengan p adalah bilangan 

prima.

• Fp adalah adalah medan berhingga dengan himpunan bilangan bulat \{0, 1, 2, …, 

p – 1\}, p bilangan prima, dan dua operasi yang didefinisikan sbb:

 1. Penjumlahan, dilakukan dalam modulus p


Jika a, b  Fp, maka a + b = r, yang dalam hal ini


r = (a + b) mod p,  0  r  p – 1

 2. Perkalian, dilakukan dalam modulus p


Jika a, b  Fp, maka a  b = s, yang dalam hal ini


s = (a  b) mod p,   0  s  p – 1

\end{document}
